\chapter{Квантовые числа бозонного дефектного поля}

Закрытие фермионной ветви спинорного накрытия оставляет безусловно выведенный объект:
поле параметра порядка $Q$ и его конечные проекторные дефекты. Теперь
необходимо определить их физический тип исключительно из действующих
представлений, общего следа и топологии.

\section{Представление поля}

Поле
\begin{equation}
 Q(x)=R(x)-\frac13I_3,
 \qquad Q^T=Q,
 \qquad \Tr Q=0
\end{equation}
не имеет пространственно-временных тензорных индексов. В текущем
ковариантном чтении оно является лоренцевым скаляром. Его пять компонент
образуют вещественное симметричное бесследовое представление внутреннего
$SO(3)_{\rm fam}$. Инфинитезимальный аудит дал
\begin{equation}
 -\sum_{a=1}^3L_a^2=6I_5,
\end{equation}
то есть внутренний спин $j=2$. Это семейный spin-two, а не
пространственно-временное поле спина два.

Калибровочная алгебра действует как
\begin{equation}
 \pi(a)=I_3\otimes\pi_{\rm SM}(a),
\end{equation}
тогда как $Q$ действует как $Q\otimes I_{15}$. Поэтому
\begin{equation}
 [Q\otimes I_{15},I_3\otimes\pi_{\rm SM}(a)]=0
\end{equation}
для всех двенадцати генераторов
$su(3)_c\oplus su(2)_L\oplus u(1)_Y$. Явный матричный остаток равен нулю.

\begin{theorem}[Первичная идентификация поля]
$Q$ преобразуется как
\begin{equation}
 (\mathbf1,\mathbf1,0)_{\rm SM}\otimes\mathbf5_{SO(3)_{\rm fam}}.
\end{equation}
Следовательно, оно является калибровочным синглетом Стандартной модели,
но нетривиальным семейным параметром порядка. Оно не является
электрослабым дублетом Хиггса.
\end{theorem}

\section{Каким образом поле может быть видно}

У spin-two представления нет инвариантного линейного функционала. Это
подтверждается нулевой размерностью общего ядра трёх генераторов и
тождеством $\Tr Q=0$. Поэтому семейно-инвариантный линейный портал к
скалярам или калибровочной плотности отсутствует.

Первый скалярный инвариант равен $\Tr Q^2$. Симметрии допускают
\begin{equation}
 \lambda_{QH}\Tr(Q^2)H^\dagger H,
 \label{eq:v6-Q-Higgs-portal}
\end{equation}
а на более высокой размерности ---
\begin{equation}
 c_{QF}\Tr(Q^2)\Tr(F_{\rm SM}^{\mu\nu}F^{\rm SM}_{\mu\nu}).
\end{equation}
Это стандартная логика портала Хиггса для калибровочно-синглетных полей; исходный
мотив см. в \path{arXiv:hep-ph/0605188}. Но коэффициенты
$\lambda_{QH}$ и $c_{QF}$ текущим родительским следом не выведены.
Следовательно, разрешённый портал ещё не является предсказанным
взаимодействием.

Прямая связь с фермионами могла бы иметь форму
\begin{equation}
 \sum_s y_s Q_{ij}\bar\psi_s^i\psi_s^j,
 \qquad s=Q,L,u,d,e.
\end{equation}
Однако физический corner-аудит уже установил пять независимых видовых
направлений и четырёхмерную центрированную неоднозначность. Обычные
внутренние флуктуации семейное поле также не создают. Поэтому ни один
универсальный ненулевой $y_s$ не выбран.

\begin{nogocriterion}[Разрешённый портал не является выведенным]
Нельзя отождествлять калибровочную нейтральность с невидимостью, но также нельзя
считать взаимодействие полученным только потому, что симметрия допускает
$\Tr(Q^2)H^\dagger H$ или семейную фермионную билинейную форму.
\end{nogocriterion}

\section{Какой топологический заряд принадлежит самому полю}

Для фиксированной базовой ориентации
\begin{equation}
 \pi_2(\mathbb{RP}^2)=\mathbb Z.
\end{equation}
Но поле зависит от директора только через
\begin{equation}
 P(n)=nn^T.
\end{equation}
Два подъёма имеют степени
\begin{equation}
 \deg(n)=+1,
 \qquad
 \deg(-n)=-1,
\end{equation}
и при этом поточечно дают один объект:
\begin{equation}
 P(n)=P(-n).
\end{equation}
На двухстах случайных направлениях максимальная разность проекторов равна
нулю.

Это точная реализация известной неоднозначности зарядов дефектов при
нетривиальном действии $\pi_1$ на высшей гомотопии; современный анализ
через накрытия дан в \path{arXiv:2109.02892}. Для самого непомеченного
бозонного поля внутренним является модуль заряда $|k|$. Знак требует
выбора подъёма, базовой ориентации или дополнительной метки спинорного накрытия.

Следовательно, классы $+15/-15$ остаются корректными ориентированными
граничными подъёмами, но не являются двумя различимыми локальными
состояниями одного только $Q$. Минимальный бозонный дефект несёт
топологическую величину $|k|=1$ и коэффициентную величину $15$ после
выбора подъёма.

\section{Наблюдаемые заряды и статистика}

Из коммутанта следуют точные нулевые квантовые числа:
\begin{equation}
 (C_2^{SU(3)},C_2^{SU(2)},Y,Q_{\rm em})=(0,0,0,0).
\end{equation}
Барионное и лептонное числа также не выведены: они не являются
топологическим зарядом поля $Q$ и требуют отдельного связывания с
фермионным сектором.

Полная Real WZW/Pfaffian-фаза предыдущего аудита равна $+1$, ориентация
линия Пфаффиана не построена, а носитель нулевых мод спинорного накрытия закрыт.
Поэтому оснований назначать дефекту фермионную статистику нет.
Классический объект построен из бозонного скалярного поля, но его строгая
квантовая статистика требует анализа фундаментальной группы пространства
конфигураций и ограничения Финкельштейна--Рубинштейна; см. DOI
\path{10.1063/1.1664510}.

\begin{proposition}[Строгий текущий статус материи]
Выведен нейтральный бозонный топологический кандидат семейного скрытого
сектора. Он не отождествлён ни с одной частицей Стандартной модели, не
имеет выведенного электрического, цветового, слабого, барионного или
лептонного заряда и пока не обладает вычисленной квантовой статистикой.
\end{proposition}

Это не отрицание материи, а первое точное ограничение её типа. Следующий
проверка должна квантовать трансляционные и ориентационные коллективные
координаты конечного дефекта, вычислить топологию его пространства
конфигураций и проверить, допускается ли нетривиальное
Финкельштейн--Рубинштейн представление без возвращения закрытого
фермионный носитель.

\begin{protocol}
\begin{enumerate}
 \item лоренцев тип $Q$: \textbf{скаляр};
 \item внутреннее семейное представление: \textbf{пятёрка, $j=2$};
 \item представление Стандартной модели: \textbf{$(1,1,0)$};
 \item поле Хиггса: \textbf{нет};
 \item линейный семейно-инвариантный портал: \textbf{отсутствует};
 \item квадратичный портал Хиггса: \textbf{разрешён, но не выведен};
 \item уникальная связь с фермионами: \textbf{не выведена};
 \item собственный заряд непомеченного $Q$-дефекта: \textbf{$|k|$};
 \item ориентированные классы $+15/-15$: \textbf{требуют подъёма};
 \item калибровочные заряды Стандартной модели: \textbf{нулевые};
 \item фермионная статистика: \textbf{не выведена};
 \item текущая идентификация: \textbf{нейтральный бозонный
 топологический кандидат};
 \item следующая проверка: \textbf{квантование коллективных координат}.
\end{enumerate}
\end{protocol}

Машинный сертификат сохранён в
\path{s2t_v6_bosonic_defect_field_identification_gate_results.json}.